% =============================================================================
% 2. INTRODUCTION
% =============================================================================
\section{Introduction}
\label{sec:intro}

Coding agents that resolve real software engineering issues have moved from
research demos to deployed systems~\citep{swebench, sweagent,
openhands}. Their progress, however, has been driven almost entirely by
scaling synthetic agent-trajectory data during post-training: pipelines such
as SWE-Gym~\citep{swegym}, r2e-gym~\citep{r2egym}, and
swe-smith~\citep{swesmith} curate or synthesise trajectories that
imitate human or LLM behaviour on issue-resolution tasks. The base model
these pipelines start from is typically a code LLM trained with next-token
prediction (and, in some cases, random-span FIM) on internet-scale
code~\citep{qwencoder, deepseekcoder, starcoder}. Between
these two stages lies a training-time gap: the base model is rarely
optimised for the conditioning structure that agentic post-training will
later demand. We treat this gap as an opportunity for a dedicated
\emph{mid-training} stage that aligns the base model with agent-relevant
inductive biases before agent-specific data is introduced.

% =============================================================================
% FIG 1 [必须]: 第一页 teaser figure
%   (left)  code snippet with call site highlighted in 4 colours, labelled:
%           caller-context / call / return / downstream-use
%   (middle) agent trajectory with the same 4 colours, labelled:
%           history h_t / action a_t / observation o_{t+1} / continuation
%   (right) bar chart: SWE-Bench-Verified across 7B/14B/32B,
%           grouped pairs (post-only vs +FIM-Midtrain)
% =============================================================================
\begin{figure}[ht]
\centering
\includegraphics[width=\linewidth]{figs/teaser.pdf}
\caption{The four-part decomposition of a function call site (left) is
structurally isomorphic to a single step of a coding agent (middle):
pre-call context $\leftrightarrow$ history, call $\leftrightarrow$ action,
return $\leftrightarrow$ external observation, downstream use
$\leftrightarrow$ continuation. Function-aware FIM mid-training exploits
this correspondence as a self-supervised prior, yielding consistent
SWE-Bench-Verified gains of $+2.8/+3.0/+3.3$ across 7B/14B/32B Qwen2.5-Coder
bases (right; solid bars = SWE-Bench-Verified, hatched bars =
SWE-Bench-Lite).}
\label{fig:teaser}
\end{figure}

Our central observation is that the inductive bias required by a coding
agent already exists in ordinary code, but in a shape that left-to-right
pretraining systematically under-exposes. At each step an agent maintains a
history $h_t$, samples an action $a_t \sim \pi(a_t \mid h_t)$, receives an
observation $o_{t+1}$ produced by an external process, and continues
conditioned on the entire trace. This four-part decomposition---context,
action, externally-computed return, continuation---is precisely the
decomposition of a function call site: pre-call code that establishes intent
and binds arguments; the call itself; a return value produced by code
outside the immediate scope; and downstream code that consumes the return
value. A model trained to reason \emph{bidirectionally} about function-level
dependencies must learn to reconstruct a callee's behaviour from caller
context and downstream usage, which is the same competence required to
predict an agent's continuation given a history and a tool return. The
correspondence is structural rather than literal---FIM training conditions
on a given suffix while agent inference generates one---but it suggests
that representations induced by the former should transfer to the latter,
an empirical question we address in Section~\ref{sec:exp}.

The fill-in-the-middle objective itself is not new: it has been part of code
LLM pretraining since~\citet{bavarian2022fim}, and recent code LLMs mix
random-span FIM into their pretraining
corpora~\citep{qwencoder, deepseekcoder, starcoder}. We
argue, however, that random-span FIM is poorly aligned with the agentic
conditioning structure for three reasons. \emph{(i) Span boundaries are
syntactically arbitrary:} most masked spans cut through expressions or
partial statements and provide only weak signal about function-level
dependencies. \emph{(ii) There is no reasoning supervision:} the model
fills the span directly, with no intermediate rationale mirroring an
agent's think-then-act pattern. \emph{(iii) The objective is dissolved
into pretraining:} by the time agent post-training begins, any structural
prior conferred by FIM has been amortised across trillions of unrelated
tokens. We address all three points. We select masking targets at function
granularity using program dependency graph analysis combined with a
complexity--inferability double criterion that is deliberately
base-model-agnostic (Section~\ref{sec:method:selection}). We embed
chain-of-thought rationales \emph{inside} the FIM middle span so that the
model learns to produce a reasoning trace consistent with the eventual
code (Section~\ref{sec:method:cot}). And we apply the objective at a
dedicated mid-training stage, concentrating its signal immediately before
agentic post-training rather than spreading it across pretraining.

We evaluate this recipe along three axes of robustness. \emph{Across model
sizes}, mid-training improves SWE-Bench-Verified at 7B
($15.0 \to 17.8$, $+2.8$), 14B ($26.2 \to 29.2$, $+3.0$), and 32B
($31.8 \to 35.1$, $+3.3$): the gain is consistent across the scales we
evaluate, indicating that the structural prior is not absorbed by larger
pretrained models and remains useful in the practical deployment range.
\emph{Across post-training pipelines}, mid-training improves both r2e-gym
and swe-smith on the same 7B base, with the swe-smith pairing improving by
$+5.3$ points on SWE-Bench-Verified. \emph{Across base-model families},
mid-training transfers from Qwen2.5-Coder to Qwen3-8B
(\textcolor{red}{$30.2 \to 33.4$ on SWE-Bench-Verified}, paired with the
swe-lego post-training pipeline), suggesting the prior is not specific to
a single pretraining recipe.

A second set of results probes our motivating hypothesis. We note first
that all benchmarks in this paper are evaluated on checkpoints that have
undergone the full pipeline---either post-training alone (baseline) or
mid-training followed by post-training (ours)---because a mid-trained
checkpoint without subsequent instruction-style post-training has degraded
instruction-following ability and is not directly comparable to
instruction-tuned baselines. Under this protocol, we observe that agentic
post-training alone substantially regresses non-target capabilities at
14B: r2e-gym alone drops LiveCodeBench from $37.2$ to $24.1$, BFCL from
$23.2$ to $15.7$, and tau-bench from $5.7$ to $0.4$. Adding mid-training
before the same post-training largely restores them to $35.2$, $18.2$, and
$4.4$ respectively---gains of $+11.1$, $+2.5$, and $+4.0$ over post-training
alone. Because the mid-training corpus contains only Python code with no
tool-use data, mid-training a priori has no reason to affect tau-bench or
BFCL; the fact that it does, and in the same direction as the in-domain
SWE-Bench gains, provides direct evidence for the isomorphism between
function calls and agent tool calls.

\paragraph{Contributions.}
\textbf{(1) Conceptual framing.} We articulate a structural correspondence
between function call sites and the action--observation--continuation loop
of coding agents, and motivate function-granularity FIM as a self-supervised
prior for agent capability that exists at internet scale in ordinary code.
\textbf{(2) Method.} We introduce a function-aware FIM mid-training pipeline
combining program dependency graph analysis, a hand-designed
complexity--inferability double criterion that is base-model-agnostic, and
chain-of-thought rationales embedded inside the FIM middle span.
\textbf{(3) Three-axis robustness.} We empirically validate the recipe along
three independent axes---model size (7B/14B/32B), post-training pipeline
(r2e-gym, swe-smith), and base-model family (Qwen2.5-Coder,
Qwen3)---observing consistent gains in every configuration, with
SWE-Bench-Verified improvements of $+2.8/+3.0/+3.3$ at 7B/14B/32B.
\textbf{(4) Direct test of the motivating hypothesis.} The same mid-training
corpus---Python code only, no tool-use data---yields gains over
post-training-only baselines on tau-bench ($+4.0$), BFCL ($+2.5$), and
LiveCodeBench ($+11.1$ recovery from post-training-induced regression),
providing direct evidence that the function-call inductive bias survives
agentic post-training and transfers across task families.
\textbf{(5) Open release.} We release the 990-repository decontaminated
corpus (78K filtered self-contained Python files, $\approx 33\mathrm{M}$
lines), the full selection pipeline, and all mid-training checkpoints to
support reproduction.